\chapter{Structure of Single-Electron Atoms}

%===========================================

%===========================================

\section{Solution of the Schrödinger Equation}

The electrostatic potential is
\[
V = -\frac{Ze^2}{4\pi\varepsilon_0 r}
\]

\subsection{Separation of Relative Motion}

Transforming the individual motions of nucleus and electron into center-of-mass motion and relative motion, we use center-of-mass position and relative position:
\[
\mathbf{X} = \frac{m_e}{M}\mathbf{x}_e + \frac{m_N}{M}\mathbf{x}_N,
\quad
\mathbf{x} = \mathbf{x}_e - \mathbf{x}_N
\]
where $M$ is the total mass and $\mu$ is the \emph{reduced mass}:
\[
M = m_e + m_N, \quad \mu = \left(\frac{1}{m_e} + \frac{1}{m_N}\right)^{-1}
\]
Expressing momenta:
\[\begin{aligned}
\mathbf{p}_e &= m_e\dot{\mathbf{x}}_e = m_e\dot{\mathbf{X}} + \frac{m_e m_N}{M}\dot{\mathbf{x}},\\
\mathbf{p}_N &= m_N\dot{\mathbf{x}}_N = m_N\dot{\mathbf{X}} - \frac{m_e m_N}{M}\dot{\mathbf{x}}
\end{aligned}\]
The kinetic energy becomes:
\[
E_k = \frac{p_e^2}{2m_e} + \frac{p_N^2}{2m_N} = \frac{1}{2}M\dot{X}^2 + \frac{1}{2}\mu\dot{x}^2 = \frac{P^2}{2M} + \frac{p^2}{2\mu}
\]
Since the relative motion wave function is unaffected by the center-of-mass kinetic energy operator, the relative Hamiltonian is:
\[
\hat{H}_{\text{rel}} = -\frac{\hbar^2}{2\mu}\nabla^2 + V
\]

\subsection{Separation of Variables}

Writing $\psi(r,\theta,\varphi) = R(r)Y(\theta,\varphi)$, we obtain:
\[
\hat{L}^2 Y = \hbar^2 l(l+1) Y\]
\[
-\frac{\hbar^2}{2\mu R}\left(r^2\frac{d^2R}{dr^2} + 2r\frac{dR}{dr}\right) + Vr^2 + \frac{\hbar^2 l(l+1)}{2\mu r^2}r^2 = Er^2
\]
The angular equation gives spherical harmonics. For the radial equation, substituting $u = rR$:
\[
-\frac{\hbar^2}{2\mu}\frac{d^2u}{dr^2} + V_{\text{eff}}u = Eu
\]
where the \emph{effective potential} is
\[
V_{\text{eff}} = -\frac{Ze^2}{4\pi\varepsilon_0 r} + \frac{\hbar^2 l(l+1)}{2\mu r^2}
\]
The second term is the \emph{centrifugal barrier}, representing the quantum mechanical version of the classical centrifugal potential $L^2/(2\mu r^2)$.\\
The solutions are:
\[
\psi_{nlm} = \sqrt{\left(\frac{2Z}{na}\right)^3\frac{(n-l-1)!}{2n[(n+l)!]}}e^{-\rho/2}\rho^l L_{n-l-1}^{2l+1}(\rho)Y_l^m(\theta,\varphi)
\]

\[
E = -\frac{Z^2\mu e^4}{8\varepsilon_0^2 h^2 n^2} , \quad \rho = \frac{2Zr}{na}, \quad a = \frac{4\pi\varepsilon_0\hbar^2}{\mu e^2}
\]
where $L_{n-l-1}^{2l+1}(x)$ are the \emph{associated Laguerre polynomials}.

\section{Properties of Atomic Orbitals}

\subsection{Shells, Subshells and Orbitals}

The electron wave function specified by principal quantum number $n$, angular quantum number $l$, and magnetic quantum number $m$ is called an \emph{atomic orbital}.\\
Orbitals with the same $n$ form a \emph{shell} (K, L, M, N... for $n=1,2,3,4...$). Orbitals with the same $n$ and $l$ form a \emph{subshell} (s, p, d, f... for $l=0,1,2,3...$).

\subsection{Radius, Nodes and Radial Distribution}\label{node}

The expectation value of orbital radius is:
\[
\langle r \rangle = \left[3n^2 - l(l+1)\right]\frac{a}{2Z}
\]
The \emph{radial distribution} represents the probability in a spherical shell:
\[
P(r)\,dr = \int_S |\psi|^2\,da\,dr = r^2|R(r)|^2\,dr
\]
(using the normalization of spherical harmonics $\int|Y|^2d\Omega = 1$).
Thus $P(r) = r^2|R|^2$.\\ \\
Setting $dP/dr = 0$ gives extrema. When $P=0$ we have \emph{radial nodes}, and when $P=P_{\text{max}}$.\\ \\
The number of nodes satisfies: radial nodes $= n-l-1$, total nodes $= n-1$.\\
The larger $n$ and smaller $l$, the wider the radius.

\subsection{Spectra and Ionization Energy of Single-Electron Atoms}

When energy difference equals photon energy, electrons transition between orbitals:
\[
\Delta E = -\frac{Z^2\mu e^4}{8\varepsilon_0^2 h^2}\left(\frac{1}{n_f^2} - \frac{1}{n_i^2}\right) = h\nu = hc\tilde{\nu}
\]
Defining the \emph{Rydberg constant} for infinite nuclear mass:
\[
R_\infty = \frac{m_e e^4}{8\varepsilon_0^2 h^3 c}
\]
For an atom with nuclear mass $m_N$, using reduced mass $\mu$:
\[
R_M = \frac{\mu}{m_e}Z^2 R_\infty
\]
Thus the spectral wavenumber and energy levels are:
\[
\tilde{\nu} = R_M\left(\frac{1}{n_f^2} - \frac{1}{n_i^2}\right), \quad E_n = -\frac{hcR_M}{n^2}
\]
(For spectral transitions, since photon has spin 1, angular momentum conservation requires $\Delta l = \pm 1$. Single-electron atoms approximately satisfy this selection rule, though it doesn't significantly affect energy level calculations.)\\ \\
The \emph{ionization energy} (first ionization potential) is the energy required to remove an electron from the ground state to infinity. For single-electron atoms, transitioning from $n=1$ to $n=\infty$:
\[
I = hcR_M
\]

\section{Fine Structure and Corrections}

\subsection{Relativistic Correction}

In relativity, time and space transformations involve the Lorentz factor $\gamma = (1-(v/c)^2)^{-1/2}$.\\
Kinetic energy and momentum are:
\[
T = (\gamma-1)mc^2, \quad p = \gamma mv
\]
Expanding the kinetic energy expression to order $p^4$:
\[
T = \sqrt{p^2c^2 + m^2c^4} - mc^2 = mc^2\left(\sqrt{1+\left(\frac{p}{mc}\right)^2} - 1\right) = \frac{p^2}{2m} - \frac{p^4}{8m^3c^2} + \cdots
\]
The first term is the classical kinetic energy; the second is the \emph{relativistic correction}. The corrected Hamiltonian is:
\[
\hat{H}' = -\frac{\hat{p}^4}{8m^3c^2}
\]
The first-order energy correction is:
\[
\Delta E = \langle\hat{H}'\rangle = -\frac{1}{8m^3c^2}\langle\psi|\hat{p}^4|\psi\rangle
\]
Since $\hat{p}$ is Hermitian, $\hat{p}^2$ is also Hermitian, so:
\[
\langle\psi|\hat{p}^4|\psi\rangle = \langle\hat{p}^2\psi|\hat{p}^2\psi\rangle
\]
From the Schrödinger equation we have:
\[
\frac{\hat{p}^2}{2m}\psi + V\psi = E\psi
\quad\Rightarrow\quad
\hat{p}^2\psi = 2m(E-V)\psi
\]
The expectation of radius powers are:
\[
\langle r^{-1}\rangle = \frac{Z}{n^2 a}, \quad \langle r^{-2}\rangle = \frac{Z^2}{n^3(l+1/2)a^2}
\]
Substituting gives:
\[
\Delta E_{\text{rel}} = -\frac{E_n^2}{2mc^2}\left(\frac{4n}{l+1/2} - 3\right) \propto Z^4
\]

\subsection{Spin-Orbit Coupling}\label{soc}

Although the electron does not physically rotate about the nucleus, it possesses angular momentum and charge, producing a magnetic moment and magnetic field. The energy splitting due to the interaction between the spin magnetic moment and the orbital magnetic field is called \emph{spin-orbit coupling}.\\ \\
From classical electrodynamics: In the electron's rest frame, the nucleus (charge $Ze$) orbits the electron, creating a magnetic field. With angular momentum $L = \mu r^2\omega$ (using reduced mass $\mu$):\\
The magnetic field from the orbiting nucleus:
\[
\mathbf{B} = \frac{\mu_0 I}{2r}, \quad I = \frac{Ze}{T} = \frac{Ze\omega}{2\pi}
\]
Using $\omega = L/\mu r^2$ and $c^2 = 1/\varepsilon_0\mu_0$:
\[
\mathbf{B} = \frac{1}{4\pi\varepsilon_0}\frac{Ze}{\mu c^2 r^3}\mathbf{L}
\]
Consider a classical image of a spinning electron, the magnetic moment and angular momentum is given by:
\[
\mathbf{\mu} = -\frac{e\pi r^2}{T}
\quad
\mathbf{S} = \frac{2\pi m_er^2}{T}
\]
We define the \emph{gyromagnetic ratio} of the electron:\label{gyro}
\[\gamma_e = \frac{\mathbf{\mu}}{\mathbf{S}} = -\frac{e}{2m_e}\]
Hence the electron's spin magnetic moment (with $g$-factor $g_e \approx 2.002$ due to relativistic effects):
\[
\mathbf{\mu}_e = g_e \gamma_e \mathbf{S} = -\frac{g_e e}{2m_e}\mathbf{S}
\]
The coupling Hamiltonian is $\hat{H} = -\mathbf{\mu}\cdot\mathbf{B}$. Accounting for the transformation from non-inertial to inertial frames (Thomas precession, which effectively reduces $g$ by 1):
\[
\hat{H}_{\text{SO}} = \frac{1}{4\pi\varepsilon_0}\frac{Ze(g_e-1)}{2\mu^2 c^2 r^3}\mathbf{S}\cdot\mathbf{L}
\]
Due to coupling, $\mathbf{L}$ and $\mathbf{S}$ no longer individually commute with $\hat{H}$. However, $\mathbf{L}^2$, $\mathbf{S}^2$, and the \emph{total angular momentum} $\mathbf{J} = \mathbf{L} + \mathbf{S}$ remain conserved. We express $\mathbf{S}\cdot\mathbf{L}$ using:
\[
\mathbf{J}^2 = (\mathbf{L}+\mathbf{S})\cdot(\mathbf{L}+\mathbf{S}) = \mathbf{L}^2 + \mathbf{S}^2 + 2\mathbf{S}\cdot\mathbf{L}
\]
Thus:
\[
\mathbf{S}\cdot\mathbf{L} = \frac{1}{2}\left[\mathbf{J}^2 - \mathbf{L}^2 - \mathbf{S}^2\right]
\]
With eigenvalue:
\[\frac{\hbar^2}{2}[j(j+1) - l(l+1) - s(s+1)]\]
and orbital radius expectation:
\[
\langle r^{-3}\rangle = \frac{Z^3}{l(l+1/2)(l+1)n^3 a^3}
\]
The spin-orbit coupling correction is:
\[
\Delta E_{\text{SO}} = \frac{E_n^2 (g_e-1)}{mc^2} \cdot \frac{n[j(j+1) - l(l+1) - s(s+1)]}{l(l+1/2)(l+1)} \propto Z^4
\]

\subsubsection{Spin-Orbit Coupling and Spin Flip}\label{spinflip}

The spin-orbit Hamiltonian can be written using ladder operators:
\[
\mathbf{L}\cdot\mathbf{S} = \hat{L}_z\hat{S}_z + \frac{1}{2}(\hat{L}_+\hat{S}_- + \hat{L}_-\hat{S}_+)
\]
Since ladder operators appear in pairs, the total angular momentum of the system is conserved. However, $\hat{S}_\pm$ changes $m_s$ by $\pm 1$, effectively flipping the electron spin.\\ \\
Therefore, $\hat{H}_{\text{SO}}$ acting on a singlet state $|S\rangle$ produces a triplet component $|T\rangle$. The transition probability (and rate) is proportional to the matrix element squared:
\[
w, P \propto |\langle T|\hat{H}_{\text{SO}}|S\rangle|^2
\]
For light atoms with the same $J$, the spin-orbit coupling is weak. For heavy atoms, $H_{\text{SO}}$ is larger, making spin-forbidden transitions between spectral terms with the same $J$ more likely (intersystem crossing), and causing mixing of states with different $L$ or $S$ (configuration interaction). See also \ref{spectral term} and \ref{socouple}.